\documentclass{article}

\usepackage{amsmath,amssymb}
\usepackage{xurl}
\usepackage[hidelinks]{hyperref}
\setlength{\emergencystretch}{2em}

% Shared commands for notes in docs/paper-gaps/.

\DeclareUnicodeCharacter{2080}{\ifmmode{}_{0}\else\textsubscript{0}\fi}
\DeclareUnicodeCharacter{2081}{\ifmmode{}_{1}\else\textsubscript{1}\fi}
\DeclareUnicodeCharacter{2082}{\ifmmode{}_{2}\else\textsubscript{2}\fi}
\DeclareUnicodeCharacter{2099}{\ifmmode{}_{n}\else\textsubscript{n}\fi}

\newcommand{\Lean}{\textsc{Lean}}
\ExplSyntaxOn
\NewDocumentCommand{\leanid}{m}{
  \ifmmode
    \textsf{\detokenize{#1}}
  \else
    \group_begin:
    \urlstyle{sf}
    \tl_set:Nn \l_tmpa_tl {#1}
    \tl_replace_all:Nnn \l_tmpa_tl {\_} {_}
    \exp_args:NV \path \l_tmpa_tl
    \group_end:
  \fi
}
\ExplSyntaxOff
\providecommand{\tr}{\operatorname{tr}}
\newcommand{\tnleanIssueBase}{https://github.com/LionSR/TNLean/issues/}
\newcommand{\tnleanPrBase}{https://github.com/LionSR/TNLean/pull/}
\newcommand{\tnleanDocBase}{https://sirui-lu.com/TNLean/docs/}
\newcommand{\ghissue}[1]{\href{\tnleanIssueBase#1}{GitHub issue \##1}}
\newcommand{\ghpr}[1]{\href{\tnleanPrBase#1}{PR \##1}}
\newcommand{\doclink}[1]{\href{\tnleanDocBase#1.html}{\path{#1}}}

% Dirac notation and the identity operator, as in blueprint/src/macros/common.tex.
% \providecommand keeps note-local definitions authoritative.
\usepackage{dsfont}
\providecommand{\ket}[1]{|#1\rangle}
\providecommand{\bra}[1]{\langle #1|}
\providecommand{\braket}[2]{\langle #1 | #2 \rangle}
\providecommand{\ketbra}[2]{|#1\rangle\!\langle #2|}
\providecommand{\Id}{\mathds{1}}
\providecommand{\one}{\mathds{1}}
\providecommand{\C}{\mathbb{C}}
\providecommand{\MN}[1]{M_{#1}(\C)}
\providecommand{\MD}{M_D(\C)}

% External referencing for paper sources.  The build helper writes sanitized
% auxiliary files under build/paper-aux/ and excludes duplicated source labels.
\usepackage{xr}

% Import a generated paper auxiliary file with a caller-supplied label prefix.
% The candidate paths support builds from docs/paper-gaps/, the repository root,
% and build/paper-aux/ itself.
\newcommand{\importpaperaux}[2]{%
  \IfFileExists{../../build/paper-aux/#2.aux}{%
    \externaldocument[#1]{../../build/paper-aux/#2}%
  }{%
    \IfFileExists{build/paper-aux/#2.aux}{%
      \externaldocument[#1]{build/paper-aux/#2}%
    }{%
      \IfFileExists{#2.aux}{%
        \externaldocument[#1]{#2}%
      }{}%
    }%
  }%
}

% General external reference macros
\newcommand{\paperref}[2]{\ref{#1-#2}}
\newcommand{\papereqref}[2]{(\ref{#1-#2})}

% CPSV16 source (1606.00608 / MPDO-22-12-17-2.tex) external document and shortcuts
\importpaperaux{cpsv16-}{1606.00608/MPDO-22-12-17-2-tnlean}
\newcommand{\cpsvref}[1]{\paperref{cpsv16}{#1}}
\newcommand{\cpsveqref}[1]{\papereqref{cpsv16}{#1}}

% Verdict marker: \gapnote{<kind>}{<status>}. Kind is the policy
% classification (clarification, local-correction, scope-restriction,
% unfaithful, false-source, open-gap); status is open, wip (elimination
% underway), resolved, or historical. Machine-read by the site index;
% prints a verdict line.
\newcommand{\gapnote}[2]{\par\noindent\textbf{Verdict:} #1 (#2).\par}


\title{Two-Projection Decomposition: Repeated-Projector Typo}
\author{}
\date{2026-08-12}

\begin{document}
\maketitle

\section{Source passage}

In the proof of Proposition~IV.5 of \cite{Cirac2017MPU}, the two-projection
lemma is used to decompose two projectors on $\mathbb C^D$. The source writes
\begin{align}
  P&=\bigoplus_i |0\rangle\langle 0|_i\oplus R,
  &P&=\bigoplus_i |v_i\rangle\langle v_i|_i\oplus S
\end{align}
at lines 760--763 of the local source file. The second projector must be $Q$,
not $P$. This follows from the preceding sentence, which introduces projectors
$P$ and $Q$, and from the subsequent identities involving $PQ$ and $QP$.

\section{Formalized interpretation}

The generic angle-block declarations
\leanid{LinearMap.IsSymmetricProjection.angleDefect_block} and
\leanid{LinearMap.IsSymmetricProjection.exists_orthonormal_angle_block}
therefore use symmetric projections $P$ and $Q$. They prove the exact local
action of both operators on a generic two-dimensional block. This is a local
correction of a typographical error and does not add a hypothesis or alter the
mathematical claim.

The declaration
\leanid{LinearMap.IsSymmetricProjection.reducedProjection} defines the reduced
projector as the orthogonal projector onto $\operatorname{ran}(PQ)$. The source
defines its projector blockwise; identifying that projector with the canonical
construction remains part of issue~\#6294. The associated theorems prove the
three absorption identities in lines 767--769, the range identities
\begin{align}
  \operatorname{ran}(\widetilde P)&=\operatorname{ran}(PQ),\\
  \operatorname{ran}(PQP)&=\operatorname{ran}(PQ),\\
  \operatorname{ran}(\widetilde P Q)&=\operatorname{ran}(\widetilde P).
\end{align}
The final range identity is the derived input to the equal-range right-factor
theorem, which gives the invertible map from line 770.

\section{Verdict}

\textbf{Verdict: resolved local fix (projector label and reduced projector).}
The formalized angle block uses $Q$ as the second projector, which resolves the
source's repeated $P$ label without changing the mathematical statement. The
reduced projector and its identities from lines 764--770 are also formalized.
The distinct global problem of assembling the generic and endpoint blocks into
an orthogonal direct sum remains open.

\bibliographystyle{alpha}
\bibliography{references}

\end{document}